\SourceLeaf{227}
formabunt triangulum sphaericum, et quidem iacebunt in casu secundo eodem ordine quo puncta \((1),(2),(3)\), in casu tertio vero ordine opposito. Denotando angulos illius trianguli simpliciter per \(L,L',L''\), atque perpendiculum in superficie sphaerica a puncto \(L''\) ad latus \(LL'\) ductum per \(p\), erit
\[
\sin p=\sin L\mathord{.}\sin LL''=\sin L'\mathord{.}\sin L'L'',\quad\text{atque}\quad\lambda L''=90^0\mp p
\]
\noindent valente signo superiori pro casu secundo, inferiori pro tertio. Hinc itaque colligimus
\[
\begin{aligned}
\pm\Delta&=\sin L\mathord{.}\sin LL'\mathord{.}\sin LL''=\sin L'\mathord{.}\sin LL'\mathord{.}\sin L'L''\\
&=\sin L''\mathord{.}\sin LL''\mathord{.}\sin L'L''
\end{aligned}
\]
\noindent Ceterum manifesto casus primus in secundo vel tertio comprehendi censeri potest, nulloque negotio perspicitur, \(\pm\Delta\) exhibere sextuplum soliditatis pyramidis inter puncta \(L,L',L''\) atque centrum sphaerae formatae. Denique hinc facillime colligitur, eandem expressionem \(\pm\tfrac{1}{6}\Delta\) generaliter exprimere soliditatem cuiusvis pyramidis inter initium coordinatarum atque puncta quorum coordinatae sunt \(x,y,z;\ x',y',z';\ x'',y'',z''\), contentae.

\Article{3}
\indent Superficies curva apud punctum \(A\) in ipsa situm curvatura continua gaudere dicitur, si directiones omnium rectarum ab \(A\) ad omnia puncta superficiei ab \(A\) infinite parum distantia ductarum infinite parum ab uno eodemque plano per \(A\) transeunte deflectuntur: hoc planum superficiem curvam in puncto \(A\) \textit{tangere} dicitur. Quodsi huic conditioni in aliquo puncto satisfieri nequit, continuitas curvaturae hic interrumpitur, uti e. g. evenit in cuspide coni. Disquisitiones praesentes ad tales superficies curvas, vel ad tales superficiei partes, restringentur, in quibus continuitas curvaturae nullibi interrumpitur. Hic tantummodo observamus, methodos, quae positioni plani tangentis determinandae inserviunt, pro punctis singularibus, in quibus continuitas curvaturae interrumpitur, vim suam perdere, et ad indeterminata perducere debere.

\Article{4}
\indent Situs plani tangentis commodissime e situ rectae ipsi in puncto \(A\) normalis cognoscitur, quae etiam ipsi superficiei curvae normalis dicitur. Directionem huius normalis per punctum \(L\) in superficie sphaerae auxiliaris repraesentabimus, atque statuemus
